\chapter{IN PARTENZA}\label{in-partenza}

\hfill\break

«Tutte le unità riferiscono di essere pronte per la partenza, capitano»,
disse Desai dalla poltrona del pilota. Alla sua destra c'era un capitano
dell'Aeronautica degli Stati Uniti che aveva pilotato F-22; è probabile
che fosse un bravo ragazzo e un pilota eccezionale, e dovevo tenerlo a
mente ogni volta che vedevo una delle nuove persone a bordo
dell'Olandese. Persone che non erano state con l'Unef, non erano mai
andate nello spazio prima, non erano state su Campo Alfa o su Paradiso,
non avevano catturato due navi nemiche e assaltato una base su un
asteroide sottoposta alla massima sorveglianza. Erano brave persone. Non
erano ancora la mia gente e, per essere davvero parte della squadra,
dovevano mettersi alla prova. Non solo ai miei occhi, ma agli occhi
dell'allegra banda di pirati originaria. E a quelli di Skippy. Non
dovevo lasciare che i pregiudizi nei confronti dei nuovi arrivati
influenzassero le mie azioni: si erano tutti guadagnati l'opportunità di
essere lì.

L'Olandese era ora piena di provviste per un lungo viaggio, avevamo cibo
a sufficienza per due anni e mezzo. Tutto e tutti erano stati
trasportati dalla superficie nelle navicelle thuranin pilotate da
Skippy. A parte Desai, che aveva pilotato la propria navicella: l'avevo
incoraggiata io a farlo, per dimostrare che era lei il nostro pilota
capo e tutti gli altri erano novellini.

Dalla superficie. Da terra. Già, stavo pensando alla Terra solo come a
un pianeta fra gli altri, e non come a casa. È probabile che fosse
meglio così, perché, a essere realistico, non mi aspettavo che qualcuno
di noi la rivedesse.

``Tutte le unità riferiscono di essere pronte.'' Ecco un altro
cambiamento rispetto ai nostri spensierati giorni pirateschi a bordo
dell'Olandese. Ora avevamo delle procedure. E manuali e liste di
controllo di cose da spuntare. Non confidavamo più soltanto sul fatto
che Skippy gestisse tutto dietro le quinte. Noi umani non cercavano
forse di capire come funzionasse la nave, ma almeno come far fare alla
nave quello che ci serviva. Se eravamo in grado di premere i pulsanti e
programmare un salto e la nave saltava all'incirca dove volevamo,
allora, anche se non avevamo idea di come funzionasse la tecnologia di
salto, andava abbastanza bene. Skippy mi disse che l'eventualità che gli
umani fossero in grado di pilotare la nave senza di lui, anche nella
maniera più rudimentale, confinava in modo inquietante con la
possibilità che quella umana diventasse una specie stellare. Se fossimo
stati una specie stellare, qualche caratteristica nella programmazione
di Skippy avrebbe dovuto impedirgli d'interagire con noi. Non era ancora
successo, e speravo che guadagnarci la definizione di ``specie
stellare'' ci richiedesse di comprendere la tecnologia e costruire le
nostre navi, cosa che era improbabile avvenisse durante tutto il corso
della mia vita, secondo Skippy.

Essere in grado di far fare all'Olandese quello che ci serviva era la
nostra unica speranza di tornare a casa, e Skippy lo capiva. Parlammo
cuore a cuore, o cuore a lattina di birra, di quello che si aspettava
quando avremmo trovato il Collettivo. Lui non aveva aspettative, poiché
i suoi ricordi erano ancora bloccati in modo frustrante; quello che
aveva erano speranze. Speranza che il Collettivo esistesse ancora nella
vaga versione che ricordava. Speranza che avrebbero comunicato con lui e
l'avrebbero accettato nella loro rete, civiltà, o qualsiasi cosa fosse.
Speranza che potessero spiegargli chi era, da dove veniva e come aveva
fatto a restare in orbita attorno a Paradiso su una nave abbandonata,
finché non era caduto da quell'orbita. Il dolore nella voce di Skippy
mentre ne parlava fece sperare anche a me che avrebbe trovato le
risposte che cercava.

Mi avvertì che il Collettivo avrebbe potuto non vedere di buon occhio il
fatto che avesse aiutato noi, creature biologiche tecnologicamente
arretrate, a prendere un'astronave; avrebbe potuto anche disabilitare la
nave. Non poteva fare promesse e lo capivo, apprezzavo la sua onestà.
Per quanto riguarda la mia di onestà, parlai con ciascuno dei nostri
volontari e, dopo aver fatto il mio discorso sugli obiettivi della
missione e sulle scarse probabilità che avevamo di tornare a casa,
nessuno si ritirò. Nessuno dei quattordici scienziati a bordo voleva
perdere l'opportunità di esplorare la galassia, e tutto il personale
militare era ansioso di assicurarsi che il wormhole restasse chiuso, per
sempre. Quando spiegai a un maggiore del corpo dei Marine, che indossava
una Stella di bronzo sull'uniforme, che il piano, se così si poteva
chiamare, era quello di vagare per la galassia finché Skippy non avesse
trovato un modo per contattare un Collettivo che avrebbe potuto non
esistere − e che inoltre non avevo idea di cosa sarebbe successo dopo −,
lui si limitò a fare un'alzata di spalle: «Cavolo, è molto più chiaro
della maggior parte delle istruzioni per la missione che ho ricevuto
come sottotenente in Iraq».

Il motivo per cui ero a bordo di questa probabile missione suicida è che
avevo promesso a Skippy di aiutarlo a localizzare il Collettivo e che
finora lui aveva fatto la sua parte per tenere fede al nostro accordo.
Tutti gli altri erano a bordo per senso del dovere, o spirito
d'avventura. Sciocco, forse, ma del tutto onorevole. Mia madre aveva
pianto come una fontana quando le avevo annunciato che, appena arrivato,
sarei tornato nello spazio per un tempo indefinito. Mio padre aveva
annuito, mi aveva stretto la mano e aveva cercato di sembrare stoico,
perché questo è l'atteggiamento di merda che gli uomini della mia
famiglia hanno sempre avuto, e io l'avevo abbracciato e avevamo pianto
entrambi. I miei genitori erano orgogliosi di me, senza avere idea di
quello che avevo fatto: tutto ciò che sapevano è che ero atterrato di
fronte a casa loro in una navicella aliena e che il dovere richiedeva
che tornassi in fretta nello spazio. Sulla griglia non c'erano
cheeseburger stavolta, perché la carne di manzo non si vedeva da un po'
nella cucina dei miei, ma c'era del pollo, un pollo che avevano allevato
nel cortile, ed era delizioso. Mia sorella, che lavorava a Boston, aveva
dovuto accontentarsi di una telefonata con il sottoscritto; mi aveva
augurato buona fortuna e mi aveva chiesto se fosse vero che i kristang
non sarebbero mai più apparsi nei nostri cieli. Sì, le avevo assicurato,
le cose erano tornate alla normalità. Se qualcuno riusciva a ricordare
cosa fosse la normalità, prima del Giorno di Colombo.

L'umanità aveva una nave da trasporto truppe kristang in orbita, ora
senza kristang e senza trappole esplosive. Una nave da trasporto che
poteva essere raggiunta solo dai vecchi razzi chimici, perché Skippy
aveva insistito sul fatto che non potevamo lasciare a terra neanche una
nave kristang o thuranin, in funzione di spintarella tecnologica. «La
tua specie è andata sulla vostra luna da sola, Joe, sono sicuro che
possono entrare in orbita senza il nostro aiuto.» Non c'era più un
kristang vivo nei crateri fumanti vicino a Hangzhou, Lione e Durango, e
la quantità di polvere sollevata dagli impattatori dei cannoni a rotaia
avrebbe fatto sì che l'umanità ricevesse in dono tramonti spettacolari
per un paio di mesi.

Abbassai lo sguardo sul mio iPad, da cui Skippy aveva cancellato il
software originale per sostituirlo con il suo. Mostrava lo stato dei
sistemi navali critici in un modo che aveva senso per me, per quanto
ridondante Skippy riteneva che fosse. Era anche pieno di tutti i tipi di
corsi che ci si aspettava seguissi in quanto ufficiale dell'esercito
degli Stati Uniti, e uno dei luogotenenti dell'esercito a bordo era
stato incaricato di aggiornarmi, oltre a svolgere il suo lavoro
regolare. Per quanto odiassi con passione i dettagli amministrativi, ero
determinato a non essere una lamentosa spina nel fianco in questo campo.
Avevo aquile d'argento sul colletto e, finché indossavo quelle,
l'esercito si aspettava che mi comportassi come un vero colonnello.
Accidenti. Alla mia età, avrei dovuto fare cose stupide e, si sperava,
imparare da esperienze dolorose. Non c'era tempo per questo, ora dovevo
rimediare alla mia mancanza di esperienza facendo affidamento sui miei
subordinati. Delegare. Avevo bisogno di abituarmi a delegare un sacco di
cose. Il tenente colonnello Chang era il mio ufficiale esecutivo, era
responsabile del quotidiano dell'equipaggio. Il maggiore Simms era ora
tornata a lavorare nella sua specialità, la logistica, e stava ancora
mettendo via con gesti frenetici la montagna di rifornimenti ammucchiati
a casaccio nelle stive di carico. Era fiduciosa che non avessimo
dimenticato nulla di vitale. Almeno un po' fiduciosa. Inoltre, con
Chang, aveva stilato una lista di compiti per l'equipaggio, che
prevedeva anche l'assegnazione di persone al lavoro nella cucina che
avevamo allestito in una delle stive di carico. Non c'erano cuochi tra
noi: avremmo tutti, me compreso, cucinato e fatto le pulizie a turno. In
quella crociera, ci sarebbero stati dei cheeseburger. Chang ordinò anche
a una squadra di lavoro di rimuovere i minuscoli letti thuranin e
sostituirli con comodi materassi a misura d'uomo. Per l'intrattenimento,
Skippy aveva scaricato tutto internet e ogni film, libro e videogioco
mai realizzato.

Il fatto che Chang, Simms e altri fossero incaricati delle operazioni
quotidiane mi lasciava libero di gestire importanti questioni
strategiche, come trattare con Skippy.

«Skippy, confermi che siamo pronti a lasciare l'orbita?»

«Eh? Ah, sì, certo, come vuoi.» Sembrava un po' irritato che le nostre
nuove procedure implicassero che noi umani dovessimo confermare le cose
da soli, invece di affidarci a lui per tutto. «Tutto a posto, colonnello
Joe. Fattore di curvatura 9 al suo segnale, o qualcosa del genere.»

«Skippy», dissi, dopo aver spento l'interfono che serviva per parlare
alla gente fuori dal ponte di comando e sussurrando nel mio microfono
perché Desai e il suo copilota e navigatore non sentissero, «ne abbiamo
parlato. Noi cavernicoli dovremo pilotare questa nave, dopo che avrai
trovato il Collettivo e ci avrai scaricato per il paradiso delle
intelligenze artificiali, o qualunque sia il posto in cui andrai.»

«No, non ne abbiamo parlato, sei stato tu a parlarne a me. Non stavi
ascoltando. Joe, non c'è quasi nessuna possibilità che voi idioti
riusciate a riportare questa nave sulla Terra. Se una qualsiasi cosa,
voglio dire, va storta con questa primitiva tecnologia thuranin, sarete
morti nello spazio senza possibilità di salvezza. Non hai idea di quanto
spesso modifichi i sistemi per farli funzionare. I tuoi uomini migliori
non capiscono nemmeno come funzionano i servizi igienici sulla
bagnarola.»

Tutto vero. I bagni thuranin non usavano acqua e in qualche modo
dividevano i rifiuti in singoli atomi di idrogeno, carbonio, ossigeno
ecc. con una procedura che i nostri fisici riconducevano all'alchimia
medievale, non alla scienza. I sistemi ambientali non solo riciclavano
l'ossigeno, ma ne creavano molecole da energia pura, contraddicendo la
famosa equazione di Einstein. Secondo i nostri scienziati, il modo in
cui venivano tracciati i salti usava la velocità della luce come una
variabile, non una costante immutabile. Skippy rispose che, sì, vedeva
come potesse sembrare così e, no, non ci avrebbe aiutato a capire come
funzionava davvero l'universo, perché tale conoscenza sarebbe stata
troppo pericolosa per le scimmie. Fu una magra consolazione quando ci
disse che neanche i maxolhx e i rindhalu avevano davvero capito come
funzionasse l'universo.

«Ti sfugge il punto, Skippy. La domanda non è se è probabile che
torniamo a casa, ma se è possibile. A bordo di questa nave ho settanta
persone bisognose di speranza, una qualsiasi speranza, che non moriranno
qui di colpo quando ci abbandonerai.»

«Non vi abbandonerò, Joe. Vi\ldots{} mmh, forse ti abbandonerò. Non
voglio. Com'è tipico delle scimmie, sei abbastanza divertente, a volte
mi sorprendi anche. Non è facile, devo dirti. Ho bisogno di farlo,
capisci? Devo sapere chi sono, da dove vengo.»

«Lo capisco e sono un soldato. Ti ho fatto una promessa la prima volta
che ci siamo incontrati e la manterrò. La mia specie è di nuovo al
sicuro ora, grazie a te. Di qualunque cosa tu abbia bisogno, te la devo,
non importa a quale prezzo. Questo non significa che autodistruggerò la
nave, a meno che non ce ne sia assoluta necessità.» Accanto all'app
Grande Pulsante Rosso sul mio zPhone, c'era ora un'app con il nome
``Boom''. Le opzioni che proponeva erano due: una distruggeva la nave in
trenta secondi, lasciandomi il tempo di cambiare idea, l'altra
distruggeva la nave immediatamente. La distruggeva, ci assicurò Skippy,
fino al livello subatomico, non lasciando alcuna possibile traccia che
l'Olandesefosse stato dirottato da umili umani. Noi saremmo morti,
mentre Skippy se ne sarebbe andato alla deriva nello spazio, emettendo
il suono metallico di un segnalatore di volo thuranin. Provai pietà per
qualunque nave l'avrebbe preso a bordo.

«Ricevuto e grazie. Confermato che tutte le unità e i sistemi sono
pronti per la partenza.»

Riaccesi l'interfono. «Pilota, segnala a Houston che siamo pronti per la
partenza, poi portaci fuori.»

«Signorsì, capitano», disse Desai con soddisfazione. «Intraprendo rotta,
ora.»

\hfill\break

Diversi giorni dopo -- più di quanti io volessi, perché mancava poco al
momento in cui il wormhole si sarebbe riavviato da solo, ma non
abbastanza secondo gli scienziati e gli ingegneri sia a bordo che a casa
-- ci fermammo vicino al punto in cui Skippy intendeva riaprire il
wormhole. Aveva voluto sfrecciare il più presto possibile in quella
direzione per avere più tempo per riavviarlo, e non vedeva alcun motivo
per ritardare. I nostri tecnici volevano che all'inizio l'Olandese se la
prendesse comoda, si riposasse e riprogrammasse tutti i sistemi critici.
Io stavo dalla parte di Skippy sulla questione; la diagnosi e il
controllo dei robot che eseguivano la manutenzione della nave era già in
esecuzione costante, e far effettuare al nostro equipaggio umano un
doppio controllo dell'analisi di Skippy appariva inutile. I nostri
migliori scienziati non avevano idea di come funzionasse un motore di
salto, lo stesso dicasi per i reattori, la gravità artificiale e in
pratica qualsiasi altro sistema thuranin. Quello che accettai, contro le
forti obiezioni dei nostri scienziati, fu un ritardo di trentasei ore.
Per lo più, acconsentii alle trentasei ore per dare a Simms e al suo
team logistico più tempo per mettere via tutta la nostra attrezzatura, e
per controllare che non avessimo dimenticato nulla di vitale, come il
ketchup. Quando l'Unef, che tecnicamente controllava ancora la missione,
mi esortò a ritardare altre ventiquattro ore, mi rifiutai e ordinai a
Desai di farci saltare. Se avessimo ritardato un altro giorno e avessimo
trovato qualcosa che non andava sulla nave, la Terra non avrebbe
comunque potuto fornire pezzi di ricambio.

Skippy riattivò il wormhole proprio come aveva predetto e l'Olandese
rimase sospesa nello spazio, vicino all'entrata stranamente luminosa del
passaggio.

«Skippy, devi lasciare che ce ne occupiamo noi», insistetti.

«No, davvero no. Se voi scimmie incasinaste una transizione nel wormhole
in futuro, sarebbe una tragedia per la nave. Se incasinaste questo, non
sarei più in grado di chiuderlo, e sarebbe una tragedia per la Terra.
Tieni a freno il tuo ego e fammi programmare la rotta da inserire. Dai,
Joe, sto riattivando un wormhole. Questo implica un sacco di variabili
che non avrete mai bisogno di affrontare di nuovo.»

Desai si voltò a guardarmi. Stavamo per entrare nel wormhole vicino alla
Terra, che era stato chiuso, un wormhole che Skippy aveva appena
riavviato. Gli schermi mostravano che c'era una pozza di luce
incandescente di fronte all'Olandese; nell'angolo dello schermo c'era un
orologio, che indicava il conto alla rovescia fino al momento in cui il
wormhole si sarebbe spostato al punto successivo. Un minuto e trentadue
secondi, e continuava a correre.

«Capitano, penso che dovremmo lasciare che sia Skippy a programmare la
rotta questa volta, non mi dispiace stare di nuovo a guardare.»

Mi morsi il labbro mentre riflettevo: «Skippy, come facciamo a sapere
che non c'è un'armata di navi thuranin dall'altra parte, impaziente di
scoprire perché questo wormhole si è spento all'improvviso?».

«Improbabile. Ma, se la cosa vi preoccupa, ho riprogrammato il
regolatore del wormhole, che ora ha una serie del tutto nuova di punti
di emergenza. Se c'è un'armata che aspetta in uno di quelli vecchi,
aspetterà per molto tempo. Per sempre.»

Resistetti alla tentazione di dire che era una cosa che avrebbe dovuto
dire prima di lasciare la Terra. «Grande. Super. Programma una rotta
attraverso il wormhole per noi.»

«Fatto.»

«Pilota», mi fermai. Questo era tutto. Era molto probabile che fosse
l'ultima volta che sarei stato nel raggio di cento anni luce dalla
Terra. Tutti i cheeseburger dell'universo non avrebbero potuto
rimediare. «Portaci dentro.»

«Sì, sì. Attivo il pilota automatico, ora.»

Il passaggio attraverso il wormhole fu la delusione che era sempre
stata. Un momento eravamo da una parte, il momento dopo eravamo a cento
o più anni luce di distanza.

«Transizione completa», riferì Desai, «e ci troviamo proprio nel punto
in cui Skippy aveva detto che ci saremmo trovati. Se questi schermi sono
precisi.»

«Lo sono», disse la nostra lattina di birra super-intelligente in tono
allegro. «Va tutto bene, Joe, nessun problema con la nave.»

«Fai quello che devi fare, Skippy», ordinai con un po' di apprensione.
Se ci avesse fregato in quel momento, non ci sarebbe stato nulla che
avrei potuto fare.

«Confermo.» Lo schermo mostrò che, dietro di noi, il pozzo di luce del
wormhole aveva lampeggiato, diciassette secondi prima. «Wormhole
disattivato.»

«Non si resetta, non si riavvia, niente del genere?» Lo schermo, che
aveva mostrato il wormhole come un punto luminoso, ora era vuoto. In un
profondo spazio interstellare, non c'era altro che l'Olandese nel raggio
di un anno luce.

«No, mai più. Il generatore del wormhole è stato spento, scollegato
dalla sua fonte di energia. Come ti ho detto, se decidi di usare il tuo
fagiolo magico, devi dargli trentotto ore per tornare online.
Quarantadue ore, per essere sicuri.»

Il fagiolo magico cui Skippy si riferiva era il modulo di controllo del
wormhole degli Anziani, che avevamo nella stiva di carico, collegato al
mio zPhone, e che poteva riattivare il wormhole per un uso una tantum.
Per permetterci di tornare a casa, senza Skippy. Una pianta di fagioli
magici che portava a casa. Avevo insistito per tenerlo con noi, Skippy
aveva ceduto a malincuore dopo una lunga discussione. Di nuovo, non si
trattava del nostro ritorno a casa, si trattava della possibilità.
Speranza. «Grazie. Eccellente. Dove si va?»

Skippy rispose: «Che ne dici di quella stella blu laggiù?».

Desai agitò il dito in direzione di numerose stelle blu sullo schermo:
«Ci sono un sacco di stelle blu. Quale?».

Skippy ridacchiò con un leggero bagliore blu: «Ha importanza?».

\hfill\break

\begin{center}
	\textbf{FINE}
\end{center}